\chapter{The MIPS Instruction Set in Depth}

With the five-instruction subset understood, we now look at the full MIPS
ISA: register conventions, the three instruction formats, addressing modes,
control flow and procedure calls.

\section{Design Principles}

Patterson and Hennessy summarize MIPS with four principles:
\begin{enumerate}
  \item \textbf{Simplicity favours regularity}: all instructions are 32 bits,
    few formats, fields in the same place.
  \item \textbf{Smaller is faster}: only 32 registers.
  \item \textbf{Make the common case fast}: immediates for small constants,
    PC-relative branches.
  \item \textbf{Good design demands good compromises}: three formats instead
    of one, to fit large constants and addresses.
\end{enumerate}

\section{Registers and Conventions}

\begin{center}
\begin{tabular}{lll}
\toprule
Number & Name & Use \\
\midrule
0 & \code{\$zero} & constant 0 \\
1 & \code{\$at} & reserved for the assembler \\
2--3 & \code{\$v0}--\code{\$v1} & return values \\
4--7 & \code{\$a0}--\code{\$a3} & arguments \\
8--15, 24--25 & \code{\$t0}--\code{\$t9} & temporaries (caller-saved) \\
16--23 & \code{\$s0}--\code{\$s7} & saved (callee-saved) \\
26--27 & \code{\$k0}--\code{\$k1} & reserved for the OS kernel \\
28 & \code{\$gp} & global pointer \\
29 & \code{\$sp} & stack pointer \\
30 & \code{\$fp} & frame pointer \\
31 & \code{\$ra} & return address \\
\bottomrule
\end{tabular}
\end{center}

\section{Instruction Formats}

\begin{center}
\begin{tabular}{l|c|c|c|c|c|c|}
\multicolumn{1}{c}{} & \multicolumn{1}{c}{31--26} & \multicolumn{1}{c}{25--21}
  & \multicolumn{1}{c}{20--16} & \multicolumn{1}{c}{15--11}
  & \multicolumn{1}{c}{10--6} & \multicolumn{1}{c}{5--0} \\
\cline{2-7}
R & op & rs & rt & rd & shamt & funct \\
\cline{2-7}
I & op & rs & rt & \multicolumn{3}{c|}{immediate (16)} \\
\cline{2-7}
J & op & \multicolumn{5}{c|}{address (26)} \\
\cline{2-7}
\end{tabular}
\end{center}

\begin{itemize}
  \item \textbf{R-format}: op = 0; the ALU operation is chosen by
    \code{funct}. \code{rs}, \code{rt} are sources, \code{rd} destination,
    \code{shamt} the shift amount for \code{sll/srl/sra}.
  \item \textbf{I-format}: arithmetic with immediates, loads/stores,
    branches. \code{rt} is the destination for \code{addi} and \code{lw},
    and a source for \code{sw} and \code{beq}.
  \item \textbf{J-format}: \code{j} and \code{jal}.
\end{itemize}

\subsection{Common instructions}

\begin{center}
\begin{tabular}{lll}
\toprule
Category & Instructions & Notes \\
\midrule
Arithmetic & \code{add, addu, sub, subu, addi, addiu} & \code{u}: no overflow trap \\
Multiply/divide & \code{mult, multu, div, divu, mfhi, mflo} & HI/LO registers \\
Logical & \code{and, or, xor, nor, andi, ori, xori} & immediates zero-extended \\
Shifts & \code{sll, srl, sra, sllv, srlv, srav} & \\
Constants & \code{lui} & load upper 16 bits \\
Compare & \code{slt, sltu, slti, sltiu} & set to 1 if less than \\
Memory & \code{lw, sw, lb, lbu, sb, lh, lhu, sh} & \\
Branch & \code{beq, bne} & PC-relative \\
Jump & \code{j, jal, jr, jalr} & \\
\bottomrule
\end{tabular}
\end{center}

\subsection{Large constants}

A 16-bit immediate cannot hold a 32-bit constant. Use \code{lui} to load the
upper half and \code{ori} to fill in the lower half:
\begin{lstlisting}[style=mips]
lui $t0, 0x1234        # $t0 = 0x12340000
ori $t0, $t0, 0x5678   # $t0 = 0x12345678
\end{lstlisting}
The pseudo-instruction \code{li \$t0, 0x12345678} expands to exactly this
(using \code{\$at} where necessary). Other pseudo-instructions: \code{move},
\code{blt}, \code{bgt}, \code{la}.

\section{Addressing Modes}

\begin{enumerate}
  \item \textbf{Register}: operand in a register (\code{add}).
  \item \textbf{Immediate}: operand in the instruction (\code{addi}).
  \item \textbf{Base/displacement}: memory at register + offset (\code{lw}).
  \item \textbf{PC-relative}: branch target $= PC + 4 + (\text{sign-ext}(imm) \ll 2)$.
    Reach: $\pm 2^{15}$ instructions.
  \item \textbf{Pseudo-direct}: jump target
    $= (PC+4)[31{:}28] \,\|\, \text{address} \,\|\, 00$. The 26-bit field
    plus two implied zero bits covers 256\,MB within the current region.
\end{enumerate}

\begin{example}[Encoding a branch]
At address 0x00400010: \code{beq \$t0, \$t1, L} where \code{L} is at
0x00400024. Offset $= (0x24 - 0x14)/4 = 4$. Encoding: op 4, rs 8, rt 9,
imm 4: \code{000100 01000 01001 0000000000000100} = \code{0x11090004}.
\end{example}

\begin{example}[Encoding a jump]
\code{j 0x00400040}: address field $=$ 0x00400040 $\gg 2$ $=$ 0x0100010.
Encoding: \code{000010} followed by 26 bits of \code{0x0100010}, giving
\code{0x08100010}.
\end{example}

\section{Control Flow}

\begin{example}[If-then-else]
\lstinline[style=c]|if (i == j) f = g + h; else f = g - h;| with
$f..j$ in \code{\$s0..\$s4}:
\begin{lstlisting}[style=mips]
      bne  $s3, $s4, Else    # branch on the opposite condition
      add  $s0, $s1, $s2
      j    Exit
Else: sub  $s0, $s1, $s2
Exit:
\end{lstlisting}
\end{example}

\begin{example}[Less-than branch]
\lstinline[style=c]|if (a < b) goto L;| with \code{a, b} in \code{\$s0, \$s1}:
\begin{lstlisting}[style=mips]
slt $t0, $s0, $s1    # $t0 = (a < b)
bne $t0, $zero, L
\end{lstlisting}
MIPS has no \code{blt} in hardware: comparison with \code{slt} plus
\code{beq}/\code{bne} against zero keeps the branch hardware simple.
\end{example}

\begin{example}[While loop over an array]
\lstinline[style=c]|while (save[i] == k) i += 1;| with \code{i} in
\code{\$s3}, \code{k} in \code{\$s5}, base of \code{save} in \code{\$s6}:
\begin{lstlisting}[style=mips]
Loop: sll  $t1, $s3, 2       # t1 = 4*i
      add  $t1, $t1, $s6     # t1 = &save[i]
      lw   $t0, 0($t1)
      bne  $t0, $s5, Exit
      addi $s3, $s3, 1
      j    Loop
Exit:
\end{lstlisting}
\end{example}

\section{Procedures}

Steps of a procedure call:
\begin{enumerate}
  \item Caller puts arguments in \code{\$a0}--\code{\$a3} (extra on the stack).
  \item \code{jal Proc}: saves $PC + 4$ in \code{\$ra} and jumps.
  \item Callee allocates a stack frame, saves registers it must preserve,
    does its work, puts results in \code{\$v0}/\code{\$v1}.
  \item Callee restores registers, frees the frame, \code{jr \$ra}.
\end{enumerate}

\subsection{The stack}

The stack lives at the top of memory and grows \emph{downward}. \code{\$sp}
points to the top (last used word). Push: decrement \code{\$sp}, then store.
Pop: load, then increment.

\begin{important}[Calling convention]
\begin{itemize}
  \item \textbf{Callee-saved}: \code{\$s0}--\code{\$s7}, \code{\$sp},
    \code{\$fp}, \code{\$ra} (the callee must restore them if it changes
    them; a non-leaf procedure must save \code{\$ra} since its own
    \code{jal} overwrites it).
  \item \textbf{Caller-saved}: \code{\$t0}--\code{\$t9}, \code{\$a0}--\code{\$a3},
    \code{\$v0}--\code{\$v1} (the callee may destroy them; the caller saves
    any it needs after the call).
\end{itemize}
\end{important}

\begin{example}[Recursive factorial]
\begin{lstlisting}[style=mips]
fact: addi $sp, $sp, -8      # make room for 2 words
      sw   $ra, 4($sp)       # save return address
      sw   $a0, 0($sp)       # save n
      slti $t0, $a0, 1       # n < 1 ?
      beq  $t0, $zero, L1
      addi $v0, $zero, 1     # base case: return 1
      addi $sp, $sp, 8       # pop (nothing changed)
      jr   $ra
L1:   addi $a0, $a0, -1
      jal  fact              # v0 = fact(n-1)
      lw   $a0, 0($sp)       # restore n
      lw   $ra, 4($sp)       # restore return address
      addi $sp, $sp, 8
      mul  $v0, $a0, $v0     # n * fact(n-1)
      jr   $ra
\end{lstlisting}
(\code{mul} is a pseudo-instruction for \code{mult} + \code{mflo}.)
\end{example}

\subsection{Memory layout}

From low to high addresses: reserved, \textbf{text} (code, starting at
0x00400000), \textbf{static data} (globals, accessed via \code{\$gp}),
\textbf{heap} (dynamic, grows up), and the \textbf{stack} (grows down from
0x7FFFFFFC). Heap and stack grow toward each other.

\section{From Source to Execution}

Compiler (C $\to$ assembly) $\to$ assembler (assembly $\to$ object file;
expands pseudo-instructions, builds symbol and relocation tables)
$\to$ linker (combines object files, resolves labels, fixes absolute
addresses) $\to$ loader (copies the executable into memory, sets up stack
and registers, jumps to the start).

\section{RISC vs CISC}

MIPS is RISC: fixed-length instructions, load/store, many registers, simple
addressing. CISC ISAs (x86) have variable-length instructions, memory
operands in arithmetic instructions and complex addressing. RISC makes the
hardware simpler and pipelining easier; modern x86 chips internally
translate instructions into RISC-like micro-operations.

\section*{Summary}
\begin{itemize}
  \item 32 registers with naming/usage conventions; \code{\$zero} is
    hardwired.
  \item R/I/J formats; opcode in bits 31--26; R-type uses \code{funct}.
  \item Five addressing modes; branches are PC-relative (word offset from
    $PC+4$); jumps are pseudo-direct.
  \item \code{lui}/\code{ori} build 32-bit constants; \code{slt} + branch
    builds comparisons.
  \item Procedures: \code{jal}/\code{jr \$ra}, arguments in
    \code{\$a}, results in \code{\$v}, stack grows down, callee saves
    \code{\$s} and \code{\$ra}.
\end{itemize}
